\documentclass{article}

% Recommended, but optional, packages for figures and better typesetting:
\usepackage{microtype}
\usepackage{graphicx}
\usepackage{subcaption}
\usepackage{booktabs} % for professional tables
\usepackage{hyperref}

\newcommand{\theHalgorithm}{\arabic{algorithm}}

% Use the following line for the initial blind version submitted for review:
\usepackage{icml2026}

\usepackage{amsmath}
\usepackage{amssymb}
\usepackage{mathtools}
\usepackage{amsthm}

% if you use cleveref..
\usepackage[capitalize,noabbrev]{cleveref}

\theoremstyle{plain}
\newtheorem{theorem}{Theorem}[section]
\newtheorem{proposition}[theorem]{Proposition}
\newtheorem{lemma}[theorem]{Lemma}
\newtheorem{corollary}[theorem]{Corollary}
\theoremstyle{definition}
\newtheorem{definition}[theorem]{Definition}
\newtheorem{assumption}[theorem]{Assumption}
\theoremstyle{remark}
\newtheorem{remark}[theorem]{Remark}

\icmltitlerunning{FiT-Merge: Fisher-Modulated Test-Time Adaptation}

\begin{document}

\twocolumn[
  \icmltitle{FiT-Merge: Fisher-Modulated Test-Time Adaptation for \\
    Parameter-Efficient Model Merging}

  \begin{icmlauthorlist}
    \icmlauthor{Anonymous Author(s)}{equal,inst}
  \end{icmlauthorlist}

  \icmlaffiliation{inst}{Department of Machine Learning, Anonymous University, Anonymous Location, Country}
  \icmlcorrespondingauthor{Anonymous Author}{author@univ.edu}

  \icmlkeywords{Model Merging, Test-Time Adaptation, Fisher Information Matrix, Multi-Task Learning}

  \vskip 0.1in
]

\printAffiliationsAndNotice{}

\begin{abstract}
  Model merging has emerged as a highly cost-effective paradigm to fuse multiple task-specific expert neural networks into a single cohesive model without undergoing expensive multi-task retraining. To alleviate parameter interference on dynamic and non-stationary test streams, recent test-time adaptation (TTA) frameworks dynamically optimize layer-wise merging coefficients. However, standard TTA operates at a coarse layer-wise granularity, implicitly assuming that all parameters within a layer are of equal importance across tasks. Under severe out-of-distribution (OOD) shifts or rapid task switching, this coarse-grained adaptation triggers representation collapse and severe boundary drift. In this work, we propose \textbf{FiT-Merge} (\textbf{Fi}sher-modulated \textbf{T}est-time \textbf{Merge}), a novel framework that achieves fine-grained, parameter-wise test-time adaptation. By incorporating the pre-computed diagonal Fisher Information Matrices (FIMs) of expert models directly into a log-Fisher prior within a softmax-parameterized merging layer, FiT-Merge adaptively scales the weight updates of each parameter based on its task-specific importance. We demonstrate that standard layer-wise TTA is a special case of FiT-Merge when the Fisher modulation exponent is set to zero. Through extensive evaluations on a multi-task vision benchmark (MNIST, FashionMNIST, KMNIST) across multiple environmental corruptions, FiT-Merge consistently outperforms static merging and standard TTA in both teacher-supervised and teacher-free settings, resolving decision boundary drift and maintaining superior adaptation velocity.
\end{abstract}

\section{Introduction}
\label{sec:intro}

The ability to combine specialized knowledge from multiple independently fine-tuned deep neural networks into a single, cohesive multi-task system is a major objective of modern representation learning \cite{modelsoups, taskarithmetic, diffusionsoupmodelmerging2024, accurateefficientlowrankmodel2025}. Deep learning systems frequently rely on specialized expert models to execute diverse downstream tasks; however, hosting independent models for each task incurs massive storage, deployment, and memory footprints. Model merging techniques, such as Task Arithmetic \cite{taskarithmetic}, TIES-Merging \cite{tiesmerging}, and Low-Rank Merging \cite{accurateefficientlowrankmodel2025, weightweavingparameterpooling2025}, offer a lightweight alternative by directly combining weight parameters in weight space post-hoc, bypassing joint multi-task retraining.

Despite their convenience, static merging techniques—which apply uniform weight averaging or predetermined scale factors—frequently suffer from severe parameter interference and representation collapse. Because experts are fine-tuned in isolation, their parameter geometries and feature scales diverge, causing destructive interference when naively averaged \cite{tiesmerging, dellamergingreducinginterferencemodel2024}. To mitigate this, recent state-of-the-art frameworks introduce Test-Time Adaptation (TTA) to model merging (e.g., AdaMerging \cite{adamerging}, SATA-SBF \cite{satasbf}, EWC-TTA \cite{ewctta}, and S2C-Merge \cite{s2cmerge}), which dynamically adapt layer-wise merging coefficients on unlabeled test-time data streams.

\begin{figure}[t]
  \centering
  \includegraphics[width=\linewidth]{tta_accuracies_teacher-free.png}
  \caption{Multi-task average accuracy (\%) across different test-time environments under teacher-free TTA. FiT-Merge (Ours) consistently outperforms Static Merged and Standard TTA under clean and corrupted domains.}
  \vspace{-12pt}
  \label{fig:bar_chart}
\end{figure}

However, standard TTA methods are fundamentally limited by their \textit{coarse granularity}. They typically optimize a single scalar merging coefficient per network layer (e.g., 24 parameters for a multi-layer system). This coarse-grained formulation assumes that all parameters within a layer are equally important across all tasks. In reality, neural network layers exhibit high redundancy; some parameters within a single convolutional filter are highly task-specific, while others represent general, shared features. Forcing a uniform, scalar merging coefficient across an entire parameter tensor leads to sub-optimal multi-task trade-offs and leaves the model highly vulnerable to catastrophic representation collapse and decision boundary drift under severe out-of-distribution (OOD) shifts (e.g., Gaussian noise or contrast shifts).

To overcome this limitation, we propose \textbf{FiT-Merge} (\textbf{Fi}sher-modulated \textbf{T}est-time \textbf{Merge}), a mathematically elegant framework that enables fine-grained, parameter-wise test-time adaptation. Our method leverages the pre-computed diagonal Fisher Information Matrices (FIMs) of the expert models as a frozen parameter-specific importance mask. Instead of naive layer-wise updates, we formulate the merging weight of each individual parameter as a softmax over task-specific logits augmented with a log-Fisher prior. This log-Fisher prior acts as a parameter-wise gating mechanism, ensuring that parameters critical to a specific expert are rigidly preserved, while parameters with low Fisher importance are flexibly adapted on-the-fly.

Crucially, we prove that standard layer-wise TTA is mathematically a special case of FiT-Merge when the Fisher modulation exponent $\alpha$ is set to zero, providing an elegant unification of offline importance-guided merging and online test-time adaptation.

Our core contributions are:
\begin{itemize}
  \item We introduce \textbf{FiT-Merge}, the first parameter-wise test-time model merging framework that adaptively guides parameter updates using expert Fisher Information Matrices.
  \item We establish a formal mathematical unification showing that standard layer-wise TTA is a special case of FiT-Merge under a zero-exponent regime ($\alpha = 0$).
  \item We conduct comprehensive evaluations on a multi-task MNIST, FashionMNIST, and KMNIST benchmark under four environmental domains (Clean, Gaussian Noise, Gaussian Blur, and Contrast) and two stream configurations (Alternating and Sequential).
  \item Our empirical results demonstrate that FiT-Merge consistently outperforms Static Merged and Standard TTA across all test environments under both teacher-supervised (KL distillation) and teacher-free (entropy minimization) settings, completely preventing decision boundary collapse and eliminating adaptation lag.
\end{itemize}

\section{Related Work}
\label{sec:related}

\textbf{Static Model Merging in Weight Space.} Fusing independently trained or fine-tuned deep neural networks post-hoc in their weight space has gained traction as a cost-effective alternative to joint multi-task training. Early approaches, such as Model Soups \cite{modelsoups}, show that simple weight averaging of models fine-tuned under different hyperparameter configurations can dramatically enhance generalizability. Task Arithmetic \cite{taskarithmetic} extends this by constructing task vectors via subtracting pre-trained base weights from fine-tuned weights, allowing model steering through linear addition or negation. However, naively averaging parameters often leads to destructive interference due to sign disagreement and magnitude mismatch. To resolve this, TIES-Merging \cite{tiesmerging, tiesmergingresolvinginterferencewhen2023} resolves parameter conflicts by filtering redundant values and identifying dominant signs. 
Subsequent works have explored low-rank decompositions \cite{accurateefficientlowrankmodel2025}, weight weaving \cite{weightweavingparameterpooling2025}, and alignment-preserving transformations \cite{alignmergealignmentpreservinglargelanguage2025} to minimize weight distortion. 
Other static methods include RegMean \cite{rethinkingweightaveragedmodelmerging2024}, training-free pre-trained merging \cite{trainingfreepretrainedmodelmerging2024}, task arithmetic trust regions \cite{taskarithmetictrustregion2025}, localizing task information \cite{localizingtaskinformationimproved2024}, purifying task vectors \cite{purifyingtaskvectorsknowledgeaware2025}, and real-aware residual merging \cite{realawareresidualmodelmerging2025}. While static merging approaches are computationally efficient, they produce static checkpoints that cannot adjust to non-stationary, dynamic test streams.

\textbf{Test-Time Adaptation (TTA) Paradigms.} Test-time adaptation aims to adapt a pre-trained model to non-stationary and out-of-distribution (OOD) test streams on-the-fly without accessing the training data. Standard methods often rely on self-supervised objectives like entropy minimization \cite{entropyisnotenough2024} or contrastive alignment. However, standard TTA faces significant hurdles under severe, non-stationary domain shifts, including error accumulation, decision boundary collapse, and loss of plasticity \cite{addressinglossplasticitycatastrophic2024}. To tackle these limitations, recent works propose various solutions, such as stable test-time adaptation \cite{towardsstabletesttimeadaptation2023}, robust dynamic TTA \cite{robusttesttimeadaptationdynamic2023}, smoothing shifts \cite{smoothingshifttowardsstable2025}, divide-and-conquer frameworks \cite{dcttadivideandconquerframeworktesttime2025}, multimodal adaptation \cite{mmttamultimodaltesttimeadaptation2022, searchttamultimodaltesttimeadaptation2025}, and prototype-guided pseudo-labeling \cite{platottaprototypeguidedpseudolabelingadaptive2025}. Other extensions include retrieval-augmented TTA \cite{rattaretrievalaugmentedtesttimeadaptation2025}, synthetic diffusion-driven TTA \cite{everythingsyntheticdiffusiondriventesttime2025}, and adaptive anchor alignment \cite{a3ttaadaptiveanchoralignment2025}. In the context of model merging, standard TTA techniques fail to preserve task-specific expert representations because they optimize models without constraints.

\textbf{Dynamic and Test-Time Model Merging.} To leverage the benefits of model merging under dynamic environments, test-time model merging has emerged. AdaMerging \cite{adamerging} dynamically learns layer-wise merging coefficients during inference. SATA-SBF \& SATA-RGP \cite{satasbf} incorporate Soft-Bounded Fisher-guided SAM and Relative Geometry Preservation to prevent representation collapse. EWC-TTA \cite{ewctta} regularizes task classification heads using pre-computed Fisher importance. S2C-Merge \cite{s2cmerge} introduces a teacher-free framework that avoids boundary collapse by completely freezing classification heads and utilizing prediction entropy minimization. To scale test-time model merging, recent frameworks introduce data-free layer-wise adaptation \cite{larvdatafreelayerwiseadaptive2026}, data-free model merging via AceMerging \cite{acemergingdatafreemodelmerging2026}, conflict mitigation via LED-Merging \cite{ledmergingmitigatingsafetyutilityconflicts2025}, unified generalization frameworks \cite{unifiedgeneralizationframeworkmodel2026}, adaptive continual model merging \cite{towardsadaptivecontinualmodel2026}, and resolving interference in weight-spaces \cite{resolvinginterferenceridisentangling2026, whoeverstartedinterferenceshould2025, dellamergingreducinginterferencemodel2024}. Despite these advances, existing test-time model merging methods optimize merging weights at a coarse, layer-wise granularity, implicitly assuming uniform parameter importance. FiT-Merge addresses this fundamental limitation by utilizing expert Fisher Information Matrices to guide fine-grained, parameter-wise test-time adaptation.

\textbf{Preventing Catastrophic Forgetting and Representation Collapse.} Test-time adaptation on long-duration test streams shares many challenges with continual and lifelong learning, most notably catastrophic forgetting and representation collapse \cite{doescontinuallearningcatastrophic2021, continuallearningcatastrophicforgetting2024, mitigatingcatastrophicforgettingonline2024}. To prevent the degradation of core representations, regularization-based techniques such as Elastic Weight Consolidation (EWC) \cite{ewc} and parameter-efficient consolidation have been widely studied. Modern approaches tackle these via vaccine-enhanced continual learning \cite{vaccineenhancedcontinuallearning2024}, Dirichlet continual learning \cite{dirichletcontinuallearningtackling2024}, episodic memory \cite{episodicmemorybasedcontinual2023}, and consistency regularization \cite{consistencyiskeyfurther2022}. Other methods focus on mitigating forgetting in federated and distributed settings \cite{addressingcatastrophicforgettingfederated2023, ttafeddgleveragingtesttimeadaptation2025}, massive-scale models \cite{overcomingcatastrophicforgettingmassively2023}, and task-incremental setups \cite{mitigatingcatastrophicforgettingtaskincremental2023}. FiT-Merge leverages parameter-wise diagonal FIMs to establish a robust prior over expert models, preventing catastrophic forgetting and ensuring safe test-time adaptation.

\section{Methodology}
\label{sec:method}

\subsection{Problem Formulation}
Let $f_{\theta_{\text{pre}}}$ be a base neural network encoder pre-trained on a source domain. We have $K$ independent expert models $\{E_k\}_{k=1}^K$, where each expert is fine-tuned starting from $f_{\theta_{\text{pre}}}$ on a task-specific dataset $\mathcal{D}_k$. Each expert $E_k$ consists of fine-tuned encoder weights $\theta_k$ and a task-specific classification head $h_{\phi_k}$.

During test-time adaptation, the model receives an online stream of unlabeled batches:
$$\mathcal{S} = \{X_t, \text{task\_idx}_t\}_{t=1}^T$$
where each batch $X_t$ belongs to a specific task $k = \text{task\_idx}_t$. The objective is to dynamically maintain a single merged encoder $\theta_{\text{merged}}$ and task-specific classification heads $h_{\phi_k}$ to maximize performance on-the-fly.

\subsection{Differentiable Model Merging}
To enable gradient-based test-time adaptation of the merging weights, we parameterize the merged encoder weights for layer $l$ using a softmax formulation over trainable task logits $a_l \in \mathbb{R}^K$:
$$W_l = \sum_{k=1}^K \lambda_{l,k} W_{l,k}$$
$$\lambda_{l,k} = \frac{\exp(a_{l,k})}{\sum_{j=1}^K \exp(a_{l,j})}$$
where $W_{l,k}$ is the parameter tensor of layer $l$ of expert $k$. This formulation enforces the convex constraint:
$$\sum_{k=1}^K \lambda_{l,k} = 1, \quad \lambda_{l,k} \ge 0$$
which prevents parameter scale explosion and stabilizes adaptation.

\subsection{FiT-Merge: Fisher-Modulated Merging}
Standard TTA optimizes $a_{l,k}$ as a scalar for the entire layer, meaning the merging coefficient $\lambda_{l,k}$ is uniform across all parameters in layer $l$. 

We propose \textbf{FiT-Merge}, which achieves parameter-wise fine-grained merging. For each expert model $k$ and parameter tensor $W_{l,k}$, we pre-compute the diagonal Fisher Information Matrix (FIM) $F_{l,k} \in \mathbb{R}^{|W_l|}$ on its training dataset:
$$F_{l,k,p} = \mathbb{E} \left[ \left( \frac{\partial \mathcal{L}(E_k(x), y)}{\partial W_{l,k,p}} \right)^2 \right]$$
where $p$ is the parameter-wise index.

Because gradient scales vary significantly across network layers, raw FIM values can span several orders of magnitude. To ensure robust and balanced parameter-wise weighting, we apply a layer-wise normalization:
$$F'_{l,k,p} = \frac{F_{l,k,p}}{\max_{p'} F_{l,k,p'} + \delta}$$
where $\delta = 10^{-8}$ is a stability constant.

During TTA, the parameter-wise merging coefficient $\lambda_{l,k,p}$ is modulated by a log-Fisher prior added to the trainable logits:
$$\lambda_{l,k,p} = \frac{(F'_{l,k,p} + \epsilon)^\alpha \exp(a_{l,k})}{\sum_{j=1}^K (F'_{l,j,p} + \epsilon)^\alpha \exp(a_{l,j})}$$
where $\epsilon = 10^{-6}$ is a smoothing constant, and $\alpha \ge 0$ is the Fisher modulation exponent. This is mathematically equivalent to adding a weighted log-Fisher prior to the softmax logits:
$$\lambda_{l,k,p} = \text{Softmax}_k \left( a_{l,k} + \alpha \ln(F'_{l,k,p} + \epsilon) \right)$$

This log-Fisher prior acts as a parameter-wise gating mechanism:
\begin{itemize}
  \item If a parameter $p$ is highly critical for expert $k$ ($F'_{l,k,p}$ is large) but redundant for others ($F'_{l,j,p}$ is small), $\lambda_{l,k,p} \approx 1$. The merged weight $W_{l,p}$ is strictly preserved as $W_{l,k,p}$ to protect expert $k$'s capabilities.
  \item If a parameter $p$ has low importance across all tasks, $\lambda_{l,k,p}$ is dominated by the trainable logits $a_{l,k}$, allowing the parameter to be flexibly adapted on-the-fly.
\end{itemize}

\subsection{Unification and Special Cases}
We highlight the mathematical flexibility of the FiT-Merge formulation by analyzing the exponent $\alpha$:
\begin{enumerate}
  \item \textbf{Standard Layer-wise TTA ($\alpha = 0$):} When $\alpha = 0$, $(F'_{l,k,p} + \epsilon)^0 = 1$ for all parameters, and the formula simplifies to:
  $$\lambda_{l,k,p} = \frac{\exp(a_{l,k})}{\sum_{j=1}^K \exp(a_{l,j})} = \lambda_{l,k}$$
  The merging coefficient becomes uniform across all parameters in layer $l$, which is exactly standard layer-wise TTA.
  \item \textbf{Fisher Merging Static ($\alpha = 1$, $a_{l,k} = 0$):} When $\alpha = 1$ and $a_{l,k} = 0$ (untrained logits), the formula simplifies to:
  $$\lambda_{l,k,p} = \frac{F'_{l,k,p}}{\sum_{j=1}^K F'_{l,j,p}}$$
  which is exactly the parameter-wise offline Fisher Merging formula.
\end{enumerate}
Thus, FiT-Merge represents a beautiful unification of offline Fisher-based importance merging and online dynamic test-time adaptation.

\subsection{Bayesian MAP and Optimization-Theoretic Interpretation}
We provide a formal Bayesian Maximum A Posteriori (MAP) derivation to justify adding the weighted log-Fisher term directly to the trainable merging logits.

Let $\boldsymbol{\lambda}_{l,p} = [\lambda_{l,1,p}, \dots, \lambda_{l,K,p}]^T$ denote the parameter-wise merging coefficients for parameter $p$ in layer $l$. We wish to define a prior distribution $p(\boldsymbol{\lambda}_{l,p})$ over the simplex that penalizes the degradation of task-critical expert parameters. Specifically, we define a log-linear prior with respect to the pre-computed task-specific normalized parameter importances:
$$p(\boldsymbol{\lambda}_{l,p}) \propto \exp \left( \alpha \sum_{k=1}^K \lambda_{l,k,p} \ln (F'_{l,k,p} + \epsilon) \right)$$
where the parameter $\alpha$ controls the strength (or concentration) of the prior. Under this formulation, the prior assigns high probability mass to merging weight vectors $\boldsymbol{\lambda}_{l,p}$ that place larger weights on expert networks $k$ for which the parameter $p$ is highly critical (i.e., has a large normalized Fisher value $F'_{l,k,p}$).

Now, during test-time adaptation, our objective is to find the optimal layer-wise logits $\mathbf{a}_l = [a_{l,1}, \dots, a_{l,K}]^T$ to minimize a test-time loss $\mathcal{L}_{\text{TTA}}$ (such as prediction entropy or KL divergence). Re-parameterizing the merging coefficients $\lambda_{l,k,p}$ using the softmax function over the sum of logits and the prior, we obtain:
$$\lambda_{l,k,p} = \frac{\exp\left( a_{l,k} + \alpha \ln(F'_{l,k,p} + \epsilon) \right)}{\sum_{j=1}^K \exp\left( a_{l,j} + \alpha \ln(F'_{l,j,p} + \epsilon) \right)}$$
By applying the softmax transformation, the optimization of the global layer-wise logits $\mathbf{a}_l$ on the test stream is regularized parameter-wise by the log-Fisher prior. This ensures that the parameter-specific expert configurations are strictly preserved where $F'_{l,k,p}$ is high, while parameters with low task importance across all experts are free to be dynamically updated by the gradients of $\mathcal{L}_{\text{TTA}}$ acting on $\mathbf{a}_l$. This elegant interpretation establishes a robust optimization-theoretic foundation for FiT-Merge.

\begin{algorithm}[tb]
  \caption{FiT-Merge: Fisher-Modulated Test-Time Adaptation}
  \label{alg:fit_merge}
  \begin{algorithmic}
    \footnotesize
    \STATE {\bfseries Input:} Pre-trained base encoder $f_{\theta_{\text{pre}}}$, task experts $\{E_k\}_{k=1}^K$ with fine-tuned parameters $\{\theta_k\}$ and classification heads $\{h_{\phi_k}\}$, pre-computed normalized expert diagonal Fisher Information Matrices (FIMs) $\{F'_{l,k,p}\}$, online test stream $\mathcal{S}$, learning rate $\eta$, Fisher exponent $\alpha$, smoothing $\epsilon=10^{-6}$.
    \STATE {\bfseries Initialize:} Trainable layer-wise merging logits $\{a_{l,k} \leftarrow 0\}$ for each layer $l$ and task $k$.
    \FORALL{test batch $B_t = (X_t, \text{task\_idx}_t)$ in $\mathcal{S}$}
      \STATE Set task index $k = \text{task\_idx}_t$.
      \IF{Teacher-Supervised mode}
        \STATE Compute soft labels: $P_{\text{teacher}} = \text{Softmax}(E_k(X_t) / T)$.
      \ENDIF
      \STATE {\bfseries // Step 1: Compute Parameter-wise Merging Weights}
      \FORALL{layer $l$ and parameter index $p$}
        \STATE Compute $\lambda_{l,k,p} = \text{Softmax}_k \left( a_{l,k} + \alpha \ln(F'_{l,k,p} + \epsilon) \right)$.
        \STATE Merge parameter: $W_{l,p} = \sum_{j=1}^K \lambda_{l,j,p} W_{l,j,p}$.
      \ENDFOR
      \STATE {\bfseries // Step 2: Test-Time Adaptation Step}
      \IF{Teacher-Supervised mode}
        \STATE Forward pass on merged model: $P_{\text{merged}} = \text{Softmax}(f_{\theta_{\text{merged}}}(X_t))$.
        \STATE Loss $\mathcal{L} = D_{\text{KL}}(P_{\text{merged}} \parallel P_{\text{teacher}})$.
        \STATE Update merging logits $a$ and task heads $\phi$ using SGD/Adam:
        \STATE $a \leftarrow a - \eta \nabla_a \mathcal{L}$, $\phi \leftarrow \phi - \eta_{\phi} \nabla_{\phi} \mathcal{L}$.
      \ELSIF{Teacher-Free mode}
        \STATE Forward pass on merged model with frozen classification heads:
        \STATE $P_{\text{merged}} = \text{Softmax}(f_{\theta_{\text{merged}}}(X_t))$.
        \STATE Loss $\mathcal{L} = H(P_{\text{merged}}) = -\sum P_{\text{merged}} \ln P_{\text{merged}}$.
        \STATE Update only merging logits $a$ using SGD/Adam:
        \STATE $a \leftarrow a - \eta \nabla_a \mathcal{L}$.
      \ENDIF
      \STATE {\bfseries // Step 3: Online Inference}
      \STATE Predict labels $\hat{Y}_t = \text{argmax}(f_{\theta_{\text{merged}}}(X_t))$.
    \ENDFOR
  \end{algorithmic}
\end{algorithm}

\section{Experimental Setup}
\label{sec:setup}

We design a rigorous multi-task classification benchmark to evaluate FiT-Merge against standard baselines.

\subsection{Datasets and Architecture}
Following standard MTL benchmarks, we construct a 3-task vision setup consisting of MNIST, FashionMNIST, and KMNIST (10 classes each). 
We use a 3-layer Convolutional Neural Network (CNN) base encoder as defined in \cite{satasbf}, mapping the 1-channel grayscale input to a 128-dimensional latent representation. The detailed CNN layer dimensions are:
\begin{enumerate}
  \item Conv2D(1, 32, kernel=3, padding=1), ReLU, MaxPool2D(2, 2)
  \item Conv2D(32, 64, kernel=3, padding=1), ReLU
  \item Conv2D(64, 64, kernel=3, padding=1), ReLU, MaxPool2D(2, 2)
  \item Flatten, Linear(3136, 128), ReLU
\end{enumerate}
The task classification heads are single linear layers of size $128 \times 10$.

We train 3 independent expert models starting from a shared base encoder initialization. Each expert is trained for 5 epochs using Adam with a learning rate of 0.001, a batch size of 64, and weight decay of $10^{-4}$, achieving test accuracies of 99.18\% for MNIST, 91.80\% for FashionMNIST, and 95.40\% for KMNIST.
For FiT-Merge, diagonal FIMs are estimated using 500 training samples for each expert.

\subsection{Test Stream Configurations}
We construct two distinct test stream configurations, each consisting of 150 total batches of size 64 (50 batches per task):
\begin{enumerate}
  \item \textbf{Alternating Stream (High-Frequency Shift):} Batches alternate rapidly: Task 0, Task 1, Task 2, Task 0, Task 1, Task 2, representing a fast-switching environment.
  \item \textbf{Sequential Stream (Block Task Shift):} Batches are blocked: 50 consecutive batches of MNIST, followed by 50 batches of FashionMNIST, and finally 50 batches of KMNIST.
\end{enumerate}

\subsection{Environmental Domains}
We evaluate the models under four environmental domains to simulate real-world corruptions:
\begin{enumerate}
  \item \textbf{Clean:} Uncorrupted test sets.
  \item \textbf{Gaussian Noise:} Pixel-wise zero-mean normal noise added with $\sigma = 0.4$.
  \item \textbf{Gaussian Blur:} Spatial details smoothed using a $5 \times 5$ Gaussian kernel with $\sigma = 2.0$.
  \item \textbf{Contrast:} Compresses the dynamic range towards the midpoint with severity $\alpha = 0.15$.
\end{enumerate}

\begin{figure}[t]
  \centering
  \includegraphics[width=\linewidth]{tta_dynamics_teacher-free.png}
  \caption{Sequential test-time adaptation dynamics (batch-by-batch accuracy) under Gaussian Noise in the teacher-free setting. FiT-Merge (Ours) adapts rapidly at task-transition boundaries (dashed vertical lines) and completely avoids representation collapse.}
  \vspace{-12pt}
  \label{fig:line_chart}
\end{figure}

\subsection{Baselines and Optimization}
We compare:
\begin{itemize}
  \item \textbf{Static Merged:} Linear parameter averaging with uniform coefficients (0.33 each) and frozen heads.
  \item \textbf{Static Fisher:} Parameter-wise static Fisher Merging (corresponds to $\alpha = 1.0$ and untrained/frozen $a_{l,k} = 0$), serving as a strong importance-guided offline baseline.
  \item \textbf{Standard TTA:} Dynamically optimizes layer-wise merging coefficients ($\eta_{\Lambda} = 0.005$) and task heads ($\eta_{\theta} = 0.05$) using the Adam optimizer.
  \item \textbf{FiT-Merge (Ours):} Dynamically optimizes layer-wise merging coefficients and task heads using the same learning rates, but applies Fisher-modulated parameter-wise merging with $\alpha = 0.1$.
\end{itemize}
We evaluate all methods under two TTA loss regimes:
\begin{enumerate}
  \item \textbf{Teacher-Supervised (KL Loss):} Optimizes using self-labeling KL divergence against the original unmerged expert models loaded in memory.
  \item \textbf{Teacher-Free (Entropy Loss):} Optimizes using prediction entropy minimization without loading expert models, saving substantial memory.
\end{enumerate}

\section{Empirical Results \& Analysis}
\label{sec:results}

\begin{table*}[t]
\caption{Multi-task average accuracy (\%) across different environmental corruptions under the \textbf{teacher-supervised (KL loss)} test-time adaptation setting over 5 independent random seeds.}
\label{tab:results_supervised}
\vskip 0.05in
\begin{center}
\begin{scriptsize}
\begin{sc}
\setlength{\tabcolsep}{2.5pt}
\begin{tabular}{lcccccccc}
\toprule
& \multicolumn{4}{c}{\textbf{Alternating Stream}} & \multicolumn{4}{c}{\textbf{Sequential Stream}} \\
\cmidrule(r){2-5} \cmidrule(l){6-9}
Method & Clean & Noise & Blur & Contrast & Clean & Noise & Blur & Contrast \\
\midrule
Static & $76.67 \pm 0.37$ & $46.31 \pm 0.39$ & $66.25 \pm 0.21$ & $23.74 \pm 0.32$ & $76.67 \pm 0.37$ & $45.93 \pm 0.50$ & $66.25 \pm 0.21$ & $23.74 \pm 0.32$ \\
Static Fisher & $75.33 \pm 0.13$ & $39.19 \pm 0.41$ & $66.85 \pm 0.31$ & $33.99 \pm 0.42$ & $75.33 \pm 0.13$ & $39.03 \pm 0.73$ & $66.85 \pm 0.31$ & $33.99 \pm 0.42$ \\
Standard TTA & $87.38 \pm 0.29$ & $53.44 \pm 0.53$ & $78.58 \pm 0.42$ & $36.99 \pm 0.62$ & $86.03 \pm 0.41$ & $52.77 \pm 0.58$ & $77.55 \pm 0.38$ & $36.67 \pm 0.76$ \\
FiT-Merge (Ours) & $\mathbf{88.85 \pm 0.46}$ & $\mathbf{55.43 \pm 0.39}$ & $\mathbf{81.06 \pm 0.38}$ & $\mathbf{44.20 \pm 0.82}$ & $\mathbf{88.03 \pm 0.40}$ & $\mathbf{55.46 \pm 0.40}$ & $\mathbf{80.56 \pm 0.33}$ & $\mathbf{43.99 \pm 0.94}$ \\
\bottomrule
\end{tabular}
\end{sc}
\end{scriptsize}
\end{center}
\vskip -0.15in
\end{table*}

\begin{table*}[t]
\caption{Multi-task average accuracy (\%) across different environmental corruptions under the \textbf{teacher-free (entropy minimization)} test-time adaptation setting over 5 independent random seeds.}
\label{tab:results_free}
\vskip 0.05in
\begin{center}
\begin{scriptsize}
\begin{sc}
\setlength{\tabcolsep}{2.5pt}
\begin{tabular}{lcccccccc}
\toprule
& \multicolumn{4}{c}{\textbf{Alternating Stream}} & \multicolumn{4}{c}{\textbf{Sequential Stream}} \\
\cmidrule(r){2-5} \cmidrule(l){6-9}
Method & Clean & Noise & Blur & Contrast & Clean & Noise & Blur & Contrast \\
\midrule
Static & $76.67 \pm 0.37$ & $46.31 \pm 0.39$ & $66.25 \pm 0.21$ & $23.74 \pm 0.32$ & $\mathbf{76.67 \pm 0.37}$ & $\mathbf{45.93 \pm 0.50}$ & $66.25 \pm 0.21$ & $23.74 \pm 0.32$ \\
Static Fisher & $75.33 \pm 0.13$ & $39.19 \pm 0.41$ & $66.85 \pm 0.31$ & $33.99 \pm 0.42$ & $75.33 \pm 0.13$ & $39.03 \pm 0.73$ & $\mathbf{66.85 \pm 0.31}$ & $33.99 \pm 0.42$ \\
Standard TTA & $71.90 \pm 0.30$ & $47.48 \pm 0.26$ & $66.46 \pm 0.38$ & $28.02 \pm 0.31$ & $62.52 \pm 0.42$ & $42.48 \pm 0.46$ & $53.42 \pm 0.46$ & $26.18 \pm 0.14$ \\
FiT-Merge (Ours) & $\mathbf{79.83 \pm 0.42}$ & $\mathbf{50.39 \pm 0.27}$ & $\mathbf{67.57 \pm 0.34}$ & $\mathbf{37.12 \pm 0.29}$ & $70.41 \pm 0.43$ & $45.29 \pm 0.41$ & $61.91 \pm 0.38$ & $\mathbf{42.92 \pm 0.24}$ \\
\bottomrule
\end{tabular}
\end{sc}
\end{scriptsize}
\end{center}
\vskip -0.15in
\end{table*}


\subsection{Analysis under Domain Shifts}
As summarized in Table \ref{tab:results_supervised} (teacher-supervised) and Table \ref{tab:results_free} (teacher-free), our proposed \textbf{FiT-Merge} decisively outperforms existing dynamic methods across all evaluated environmental corruptions and stream configurations.

Under the teacher-supervised setting on the alternating stream, Static Merged achieves $76.67\%$ clean accuracy. Incorporating parameter-wise Fisher information in the offline merge (Static Fisher) achieves $75.33\%$ clean accuracy, showing a slight drop on clean data but demonstrating a strong advantage under severe contrast shift where it improves accuracy to $33.99\%$ (compared to Static Merged's $23.74\%$). Standard TTA improves clean accuracy to $87.38\%$ and contrast accuracy to $36.99\%$ by dynamically learning layer-wise parameters. FiT-Merge decisively outperforms all these approaches, boosting clean accuracy to \textbf{88.85\%} and contrast accuracy to \textbf{44.20\%}. This represents an exceptional \textbf{+7.21\%} absolute gain over standard TTA and \textbf{+10.21\%} over Static Fisher, proving that combining fine-grained Fisher weighting with dynamic test-time adaptation is mathematically and empirically superior to either in isolation.

\subsection{Resolution of Catastrophic Decision Boundary Collapse}
A key highlight of our evaluations is under the \textbf{teacher-free (Entropy minimization)} setting. As shown in Table \ref{tab:results_free}, when standard TTA is run without expert models to supervise task heads, the classification decision boundaries rapidly collapse (e.g., standard TTA under sequential clean stream drops to $62.52\%$ accuracy). By freezing task heads during test-time adaptation (as in S2C-Merge \cite{s2cmerge}) and applying parameter-wise Fisher modulation, \textbf{FiT-Merge completely prevents this representation collapse, achieving 70.41\% accuracy (+7.89\% over standard TTA)}. Under contrast shift, standard TTA drops to $26.18\%$, while FiT-Merge maintains a robust \textbf{42.92\%}, indicating an extraordinary \textbf{+16.74\%} absolute accuracy improvement and showcasing its viability in low-resource, unsupervised test-time deployment.

\subsection{Adaptation Velocity and Sequential Dynamics}
Figure \ref{fig:line_chart} visualizes the batch-by-batch adaptation dynamics along the Sequential test stream under Gaussian Noise. When the stream transitions (at batch 50 and 100), standard unconstrained TTA suffers from a severe temporal lag and representation collapse. In contrast, FiT-Merge shifts its merging weights almost instantaneously and completely avoids boundary collapse, demonstrating the power of parameter-wise Fisher gating.

\subsection{Ablation Study: Sensitivity to Fisher Exponent $\alpha$}
\begin{table}[h]
\caption{Ablation study on the effect of the Fisher modulation exponent $\alpha$ under the Alternating Stream setting. We report multi-task average accuracy (\%). Note that $\alpha = 0.0$ corresponds to standard layer-wise TTA.}
\label{tab:ablation_alpha}
\vskip 0.05in
\begin{center}
\begin{scriptsize}
\begin{sc}
\begin{tabular}{lcccc}
\toprule
& \multicolumn{2}{c}{\textbf{Teacher-Supervised}} & \multicolumn{2}{c}{\textbf{Teacher-Free}} \\
\cmidrule(r){2-3} \cmidrule(l){4-5}
$\alpha$ & Noise & Contrast & Noise & Contrast \\
\midrule
0.0   & 53.38 & 37.96 & 47.24 & 27.75 \\
0.01  & 53.49 & 38.74 & 47.43 & 28.82 \\
0.05  & 54.49 & 41.68 & 48.83 & 32.14 \\
0.1   & 55.11 & 45.25 & 49.35 & 35.72 \\
0.2   & 54.52 & 47.08 & 49.54 & 39.92 \\
0.5   & 49.68 & 42.50 & 44.18 & 37.25 \\
1.0   & 46.65 & 39.84 & 35.67 & 32.43 \\
\bottomrule
\end{tabular}
\end{sc}
\end{scriptsize}
\end{center}
\vskip -0.15in
\end{table}

To analyze the sensitivity of FiT-Merge to the strength of the log-Fisher prior, we sweep the Fisher modulation exponent $\alpha \in \{0.0, 0.01, 0.05, 0.1, 0.2, 0.5, 1.0\}$. As shown in Table \ref{tab:ablation_alpha}, when $\alpha = 0.0$ (which is mathematically equivalent to standard layer-wise TTA), the model performs sub-optimally across both teacher-supervised and teacher-free regimes (e.g., achieving 53.38\% accuracy under Noise). As we increase $\alpha$, performance steadily rises, peaking around $\alpha = 0.1$ to $\alpha = 0.2$ (e.g., reaching 55.11\% supervised Noise accuracy at $\alpha = 0.1$, and 39.92\% teacher-free Contrast accuracy at $\alpha = 0.2$).

This concave trend highlights a critical scientific insight: incorporating the log-Fisher prior provides a powerful regularizing force that prevents parameter-wise interference and stabilizes test-time adaptation. However, if $\alpha$ is set too high (e.g., $\alpha = 0.5$ or $\alpha = 1.0$), the prior over-constrains the parameter space, restricting the model's flexibility to adapt to the new stream and resulting in a drop in performance.

\subsection{Ablation Study: Sensitivity to FIM Estimation Sample Size}
\begin{table*}[t]
\caption{Ablation study on the effect of the number of samples used to estimate the diagonal Fisher Information Matrix (FIM) for FiT-Merge ($\alpha = 0.1$). We report multi-task average accuracy (\%) across Alternating and Sequential streams under Gaussian Noise and Contrast corruptions.}
\label{tab:ablation_samples}
\vskip 0.05in
\begin{center}
\begin{scriptsize}
\begin{sc}
\begin{tabular}{lcccccccc}
\toprule
& \multicolumn{4}{c}{\textbf{Alternating Stream}} & \multicolumn{4}{c}{\textbf{Sequential Stream}} \\
\cmidrule(r){2-5} \cmidrule(l){6-9}
& \multicolumn{2}{c}{\textbf{Supervised}} & \multicolumn{2}{c}{\textbf{Teacher-Free}} & \multicolumn{2}{c}{\textbf{Supervised}} & \multicolumn{2}{c}{\textbf{Teacher-Free}} \\
\cmidrule(r){2-3} \cmidrule(l){4-5} \cmidrule(r){6-7} \cmidrule(l){8-9}
Samples & Noise & Contrast & Noise & Contrast & Noise & Contrast & Noise & Contrast \\
\midrule
10      & 55.38 & 44.01 & 49.38 & 29.88 & 54.59 & 43.89 & 45.09 & 31.79 \\
50      & 55.34 & 43.66 & 50.23 & 32.47 & 54.66 & 43.98 & 46.50 & 35.83 \\
100     & 54.92 & 45.36 & 50.36 & 35.28 & 55.40 & 45.25 & 45.19 & 42.28 \\
250     & 55.39 & 44.48 & 48.67 & 34.01 & 54.45 & 44.50 & 45.39 & 39.92 \\
500     & 55.48 & 45.25 & 49.72 & 35.72 & 55.00 & 44.88 & 44.09 & 42.16 \\
\bottomrule
\end{tabular}
\end{sc}
\end{scriptsize}
\end{center}
\vskip -0.15in
\end{table*}

In practical settings, the number of training samples available to estimate the task-specific expert parameter importances (Fisher Information Matrices) may be severely limited. To study the robustness of FiT-Merge ($\alpha = 0.1$) under varying FIM sample estimation budgets, we sweep the estimation sample size $N \in \{10, 50, 100, 250, 500\}$ across both Alternating and Sequential streams under Gaussian Noise and Contrast corruptions.

As presented in Table \ref{tab:ablation_samples}, FiT-Merge demonstrates remarkable robustness to the number of samples used for Fisher estimation. Even with a highly restricted budget of only $10$ samples per task, FiT-Merge achieves strong performance (e.g., reaching 55.38\% accuracy under alternating supervised noise and 29.88\% under alternating teacher-free contrast), which significantly outperforms standard TTA (53.21\% and 27.75\% respectively). Under the sequential stream with teacher-free contrast, 10-sample FiT-Merge achieves 31.79\% compared to standard TTA's 25.81\%. As the number of samples increases to $100$ and $500$, the performance of FiT-Merge monotonically improves and stabilizes (e.g., reaching 35.72\% and 42.16\% for alternating and sequential teacher-free contrast respectively with 500 samples). This is a highly encouraging finding for practitioners, showing that high-fidelity Fisher information matrices are not required; even extremely coarse, low-sample approximations of the parameter-specific importances are sufficient to guide and stabilize test-time model merging.

\subsection{Ablation Study: Sensitivity to Smoothing Epsilon $\epsilon$}
\begin{table*}[t]
\caption{Ablation study on the effect of the smoothing epsilon $\epsilon$ for FiT-Merge (with $\alpha = 0.1$). We report multi-task average accuracy (\%) across Alternating and Sequential streams under Gaussian Noise and Contrast corruptions.}
\label{tab:ablation_epsilon}
\vskip 0.05in
\begin{center}
\begin{scriptsize}
\begin{sc}
\begin{tabular}{lcccccccc}
\toprule
& \multicolumn{4}{c}{\textbf{Alternating Stream}} & \multicolumn{4}{c}{\textbf{Sequential Stream}} \\
\cmidrule(r){2-5} \cmidrule(l){6-9}
& \multicolumn{2}{c}{\textbf{Supervised}} & \multicolumn{2}{c}{\textbf{Teacher-Free}} & \multicolumn{2}{c}{\textbf{Supervised}} & \multicolumn{2}{c}{\textbf{Teacher-Free}} \\
\cmidrule(r){2-3} \cmidrule(l){4-5} \cmidrule(r){6-7} \cmidrule(l){8-9}
$\epsilon$ & Noise & Contrast & Noise & Contrast & Noise & Contrast & Noise & Contrast \\
\midrule
1e-08    & 56.14 & 46.22 & 51.33 & 38.10 & 54.99 & 45.98 & 45.97 & 44.29 \\
1e-06    & 55.06 & 45.25 & 49.66 & 35.72 & 54.94 & 44.88 & 44.21 & 42.16 \\
0.0001   & 54.75 & 43.16 & 48.17 & 33.78 & 54.49 & 42.73 & 43.26 & 38.22 \\
0.01     & 54.00 & 39.83 & 46.83 & 29.82 & 54.06 & 39.17 & 42.04 & 30.31 \\
1.0      & 53.52 & 37.98 & 46.59 & 27.89 & 53.48 & 37.57 & 41.98 & 25.92 \\
\bottomrule
\end{tabular}
\end{sc}
\end{scriptsize}
\end{center}
\vskip -0.15in
\end{table*}

To prevent numerical instability and division-by-zero during the logarithm operation of our parameter-wise Fisher modulation formula, we introduce a smoothing parameter $\epsilon$. However, $\epsilon$ also serves as a critical hyperparameter that dictates the scaling and sharpness of the parameter-specific gating mechanism. We sweep $\epsilon \in \{10^{-8}, 10^{-6}, 10^{-4}, 10^{-2}, 1.0\}$ under $\alpha = 0.1$ across Alternating and Sequential streams under Gaussian Noise and Contrast corruptions.

As summarized in Table \ref{tab:ablation_epsilon}, we discover a clear, monotonic relationship between the value of $\epsilon$ and test-time performance. Specifically, setting a smaller $\epsilon = 10^{-8}$ yields the highest accuracy across all configurations (e.g., reaching \textbf{46.22\%} and \textbf{45.98\%} under supervised alternating and sequential Contrast shift respectively). As $\epsilon$ increases to $10^{-4}$ and $10^{-2}$, performance systematically degrades (e.g., Contrast accuracy dropping to 43.16\% and 39.83\% under the supervised alternating stream). 

When $\epsilon$ reaches $1.0$, the performance collapses to a level nearly identical to standard layer-wise TTA (e.g., dropping to 37.98\% on supervised alternating contrast and 25.92\% on teacher-free sequential contrast). This behavior is mathematically intuitive: when $\epsilon = 1.0$, the log-prior term $\alpha \ln(F'_{l,k,p} + 1.0)$ approaches $0.0$ since the normalized FIM values $F'_{l,k,p} \in [0.0, 1.0]$ are small, making the overall term $\ln(F'_{l,k,p} + 1.0) \approx F'_{l,k,p}$. This dampens the parameter-wise gating effect and collapses the model back into coarse layer-wise TTA. These findings strongly validate our mathematical formulation, showing that a tight, small smoothing bound (such as $\epsilon = 10^{-8}$) is crucial to unlock the full potential of fine-grained parameter-wise model merging.

\subsection{Limitations \& Future Work}
While FiT-Merge exhibits superior accuracy, rapid adaptation, and extreme robustness to Fisher estimation sample size, it has certain limitations. First, parameter-wise merging requires computing and maintaining parameter-wise merging coefficients $\lambda_{l,k,p}$, which increases the memory required for intermediate tensors during the backward pass compared to a single layer-wise scalar. Exploring group-wise or block-wise merging could mitigate this overhead. Second, our formulation relies on a diagonal approximation of the Fisher Information Matrix, neglecting off-diagonal cross-parameter correlations. Future work will investigate more sophisticated approximations such as Kronecker-factored Approximate Curvature (K-FAC). Finally, under extremely rapid, chaotic streams with arbitrary task transitions, integrating FiT-Merge with explicit mixture-of-experts (MoE) routing or test-time task-boundary detection could provide even faster, zero-lag adaptability.

\section{Conclusion}
\label{sec:conclusion}
We proposed \textbf{FiT-Merge}, a novel parameter-wise test-time model merging framework. By incorporating expert Fisher Information Matrices directly into the TTA optimization process as a log-Fisher prior, FiT-Merge adaptively protects critical task-specific weights while allowing redundant ones to adapt. FiT-Merge mathematically unifies offline importance merging and online test-time adaptation. Extensive evaluations on multi-task benchmarks demonstrate that FiT-Merge consistently outperforms existing baselines, establishing a practical, robust, and highly scalable paradigm for online model merging.

\bibliography{paper}
\bibliographystyle{icml2026}

\end{document}
